\begin{abstract}
Large-scale unsupervised language models (LMs) acquire extensive world knowledge and demonstrate some capacity for reasoning; nevertheless, their behavior remains challenging to control precisely due to the inherently unsupervised training paradigm. Current techniques that enhance model steerability depend on gathering human annotations that indicate preferences between generated outputs, subsequently fine-tuning the LMs to align with these annotated preferences, frequently employing reinforcement learning from human feedback (RLHF). Yet, RLHF involves a multi-stage, often unstable process: initially fitting a reward model to capture human preferences, then using reinforcement learning to update the unsupervised LM, maximizing this surrogate reward while constraining deviation from the pre-trained model. This paper presents a novel reward model parameterization within the RLHF framework, making it possible to analytically derive the optimal policy associated with the reward model in closed form. This enables the standard RLHF objective to be addressed via a straightforward classification loss. The proposed method, termed \textit{Direct Preference Optimization} (DPO), yields a stable, efficient, and lightweight approach, removing the necessity for generation sampling during fine-tuning and minimizing hyperparameter tuning efforts. Empirically, DPO aligns LMs with human preferences at least as effectively as state-of-the-art approaches. In particular, DPO surpasses PPO-based RLHF in sentiment control tasks, and achieves parity or improvements in summarization and single-turn dialogue generation, all while being much easier to implement and train.
\end{abstract}

\section{Introduction}
Large-scale unsupervised language models (LMs) trained on extensive datasets have demonstrated remarkable and unexpected abilities~\citep{chowdhery2022palm, brown2020language, touvron2023llama,bubeck2023sparks}. Despite this, the data these models are trained on originates from humans who exhibit diverse objectives, competencies, and motivations. Not all such behaviors are desirable to replicate; for instance, it is advantageous for an AI coding assistant to \textit{recognize} prevalent programming errors so it can rectify them, but when it produces code, our preference is for the model to reflect the (sometimes uncommon) examples of expert programming found within its training corpus. Similarly, it is beneficial for the model to be \textit{informed} about widely held misconceptions (even if believed by 50\% of people), but undesirable for the model to endorse those misconceptions half of the time when queried. Thus, it becomes essential to differentiate the model's \emph{target behaviors and responses} from its expansive repertoire of \textit{knowledge and capabilities} to ensure the development of AI systems that are reliable, effective, and amenable to control \citep{ouyang2022training}. Although current methodologies primarily rely on reinforcement learning (RL) to align LMs with human judgments, our analysis demonstrates that the objective underpinning these RL-based methods can, in fact, be solved directly through a straightforward binary cross-entropy formulation. This realization allows for substantial streamlining of the preference optimization workflow.

In broad terms, current approaches infuse desired behaviors into language models by leveraging curated datasets of human preference judgments that exemplify responses considered safe and beneficial. The preference alignment phase typically follows an initial unsupervised pre-training phase on a massive corpus of text. The most direct technique for learning preferences is supervised fine-tuning using high-quality human-authored responses, but reinforcement learning from human (or AI) feedback (RLHF/RLAIF; \citep{christiano2017deep,bai2022constitutional}) has emerged as the most effective paradigm. RLHF approaches involve learning a reward function from human preference data, which is subsequently used in RL to adjust the language model policy towards favoring highly rewarded outputs, all while preventing substantial divergence from the original model. Though RLHF has led to models exhibiting impressive capabilities in conversation and coding, its workflow is notably more intricate than that of supervised learning. The RLHF process necessitates training multiple models and continuously sampling from the evolving policy during learning, resulting in considerable computational expense.

This work demonstrates a method for directly aligning language model behavior with human preferences, circumventing the need for explicit reward modeling or reinforcement learning techniques. We introduce \textit{{Direct Preference Optimization} (DPO)}, an approach that effectively targets the same reward maximization objective under a KL-divergence regularization as standard RLHF frameworks, but offers a much simpler implementation and training procedure. At a high level, DPO functions by boosting the log odds of preferred responses relative to those deemed less desirable, while simultaneously employing a per-example importance weighting mechanism. This mechanism is critical, as it mitigates the risk of model degradation seen when using basic probability ratio-based objectives. As with prior algorithms, DPO is grounded in a theoretical preference framework (e.g., the Bradley-Terry model; \cite{bradley1952rankanalysis}), which quantifies the compatibility between a reward function and observed preference annotations. Distinct from earlier techniques—which leverage the preference model to derive a training signal for a reward model, followed by policy optimization with respect to the reward—DPO reformulates the preference loss directly as a function of the policy through a variable substitution. Consequently, with access to human-annotated preference pairs over model outputs, DPO can train policies using a straightforward binary cross entropy loss, yielding an optimal policy with respect to an implicit reward inferred from the empirical data.

Our principal innovation is Direct Preference Optimization (DPO), which presents a reward-model and RL-free strategy for preference-based language model training. Empirically, we find that DPO matches or surpasses the effectiveness of traditional methods, including PPO-based RLHF, across a variety of preference learning scenarios such as sentiment conditioning, summarization, and dialogue modeling with models scaling up to 6B parameters.

\section{Related Work}

Progressively larger self-supervised language models are capable of addressing certain tasks in a zero-shot manner \citep{radford2019language} or with only minimal prompting \citep{gpt3,megatron,chowdhery2022palm}. Nonetheless, adapting these models to downstream tasks and ensuring that they act in accordance with user intent can be substantially enhanced through fine-tuning on collections of instructions paired with human-authored responses \citep{mishra-etal-2022-cross,sanh2022multitask,chung2022scaling,thoppilan2022lamda}. The process of `instruction-tuning' allows large language models (LLMs) to better extrapolate to novel instructions beyond those found in the training set, resulting in broader utility \citep{chung2022scaling}. Although instruction-tuning has proven effective, acquiring \textit{comparative} human assessments of model outputs is typically more feasible than gathering detailed expert demonstrations. Accordingly, subsequent research has refined LLMs using datasets consisting of human preferences, which has led to marked gains in domains such as translation \citep{kreutzer-etal-2018-reliability}, summarization \citep{stiennon2022learning,ziegler2020finetuning}, narrative generation \citep{ziegler2020finetuning}, and following instructions \citep{ouyang2022training,ramamurthy2023is}. These approaches generally involve first training a neural reward model aligned to preference data, often using a statistical preference framework like the Bradley-Terry model \citep{bradley1952rankanalysis}. Subsequently, the language model is fine-tuned to maximize rewards under this model, utilizing reinforcement learning techniques such as REINFORCE \citep{williams1992reinforce}, proximal policy optimization (PPO; \cite{schulman2017proximal}), or their derivatives \citep{ramamurthy2023is}. Another related line of investigation utilizes LLMs, already tuned for instruction adherence and human feedback, to produce synthetic preference datasets targeting characteristics like safety or benignity \citep{bai2022constitutional}; here, supervision from humans is limited to providing a textual rubric for the LLM’s assessments. Collectively, these approaches unify advances from two distinct streams of research: reinforcement learning-based language model training for various end-goals~\citep{Ranzato2015SequenceLT,paulus2018a,wu2018learning}, and foundational strategies for preference-based learning from human input \citep{christiano2017deep,kupcsik2018learning}. Despite the intuitive advantages of utilizing relative human preferences, reinforcement learning-based fine-tuning of large language models continues to pose significant practical difficulties; the present work introduces a theoretically-sound method to optimize relative preferences without relying on RL.

Beyond the realm of language-based tasks, the study of learning policies from preference signals has been explored within both the bandit and reinforcement learning literature, leading to the development of multiple methodologies. The framework known as the contextual dueling bandit (CDB; \cite{yue2012karmed,dudik2015contextual}) addresses the contextual bandit problem by utilizing preferences or ordinal comparisons between actions, in place of explicit reward feedback. Lacking absolute reward values, CDBs are analyzed through the concept of a \textit{von Neumann winner}, that is, a policy that is expected to win at least half of its encounters against every alternative policy \citep{dudik2015contextual}. In the CDB model, preference signals are acquired interactively in an online setting; this contrasts with learning from human preferences, where models are typically trained using a fixed, offline dataset consisting of preference-labelled pairs of actions \citep{yan2022human}. In a related area, \textit{preference-based RL} (PbRL) approaches address the challenge of learning from pairwise preferences sourced from an opaque `scoring' function, instead of from explicit rewards \citep{BusaFekete2014,ruiz2023dueling}. PbRL encompasses a variety of algorithms, including those designed to exploit off-policy preference information; however, these methods generally proceed by first inferring the underlying scoring (reward) function before subsequently performing policy optimization \citep{jain2013learning,BusaFekete2014,christiano2017deep,sadigh2017active,kupcsik2018learning}. Departing from this standard two-step methodology, we introduce an approach in which policy learning occurs in a single stage, with policies being directly optimized to satisfy the expressed preferences.

\section{Preliminaries}\label{section:prelims}

Here we outline the RLHF pipeline introduced by \citeauthor{ziegler2020finetuning}, and further developed by works such as \citep{stiennon2022learning, bai2022training, ouyang2022training}. This pipeline is typically structured into three major stages: (1) supervised fine-tuning (SFT); (2) preference collection and reward modeling; and (3) reinforcement learning-based optimization.

\textbf{SFT}: The RLHF process customarily commences with supervised fine-tuning of a pre-existing language model on carefully curated, high-quality examples relevant to the targeted applications (such as dialogue or summarization), yielding a model denoted $\pi^\text{SFT}$.

\textbf{Reward Modelling Phase}: During the reward modeling stage, the SFT model receives prompts $x$ and is tasked with generating pairs of candidate responses $(y_1, y_2)\sim \pi^\text{SFT}(y \mid x)$. These output pairs are then evaluated by human annotators, who are asked to indicate their preferred response, which is represented as $y_w\succ y_l \mid x$; here, $y_w$ and $y_l$ are the preferred and non-preferred options within the generated pair $(y_1, y_2)$. The observed preference data is assumed to reflect the judgments of an implicit, unobservable reward function $r^*(y, x)$. Various probabilistic models have been adopted to capture this preference process, with the Bradley-Terry (BT) \cite{bradley1952rankanalysis} model being among the most widely utilized; more general ranking models, such as Plackett-Luce \citep{plackett1975analysis, luce2012individual}, can also be employed when multiple ranked responses are available. According to the BT formulation, the underlying human preference distribution $p^*$ takes the following form:

\begin{equation}\label{eq:bradley-terry}
    p^*(y_1\succ y_2 \mid x)=\frac{\exp\left(r^*(x, y_1)\right)}{\exp\left(r^*(x, y_1)\right) + \exp\left(r^*(x, y_2)\right)}.
\end{equation}

Suppose we are given a fixed dataset of comparisons $\mathcal{D}=\bigl\{x^{(i)}, y_w^{(i)}, y_l^{(i)}\bigr\}_{i=1}^N$ sampled according to $p^*$. We introduce a parameterized reward function $r_{\phi}(x, y)$, and estimate its parameters via maximum likelihood estimation. By posing the task as binary classification, we can write the negative log-likelihood loss as:

\begin{equation}\label{eq:reward_model}
    \mathcal{L}_R(r_{\phi}, \mathcal{D}) = -\mathbb{E}_{(x, y_w, y_l)\sim \mathcal{D}}\bigl[\log \sigma(r_{\phi}(x, y_w)- r_{\phi}(x, y_l))\bigr]
\end{equation}

with $\sigma$ denoting the logistic sigmoid function. In large language model applications, the reward estimator $r_{\phi}(x, y)$ is commonly initialized using parameters from the SFT model $\pi^\text{SFT}(y \mid x)$, augmented with a linear output layer attached to the last transformer block to generate a scalar reward value \cite{ziegler2020finetuning}. To mitigate reward variance, earlier studies normalize the outputs so that $\mathbb{E}_{x,y\sim \mathcal{D}}\left[r_\phi(x, y)\right] = 0$ holds for each $x$.

\textbf{RL Fine-Tuning Phase}: In the reinforcement learning phase, the trained reward model provides feedback signals to guide the language model's updates. Consistent with earlier research~\citep{jaques2017sequence, jaques2020human}, the optimization can be expressed as

\begin{equation}\label{eq:RL}
\max_{\pi_{\theta}}  \mathbb{E}_{x\sim \mathcal{D}, y\sim \pi_{\theta}(y \mid x)}\bigl[r_{\phi}(x, y)\bigr] - \beta\mathbb{D}_{\textrm{KL}}\bigl[\pi_{\theta}(y\mid x)\mid \mid \pi_\text{ref}(y\mid x)\bigr],
\end{equation}

where $\beta$ modulates the penalty relative to a reference policy $\pi_\text{ref}$, which typically corresponds to the initial SFT model $\pi^\text{SFT}$. 
In practical setups, the policy $\pi_\theta$ is initialized from $\pi^\text{SFT}$. The KL regularization term acts as a safeguard against excessive divergence from the regions where the reward model is reliable, ensuring output diversity and curbing the risk of collapsing to limited, high-reward outputs. Since natural language output is discrete, this formulation is not directly differentiable, necessitating reinforcement learning algorithms for optimization. The standard protocol~\citep{ziegler2020finetuning, stiennon2022learning, bai2022training, ouyang2022training} forms a reward surrogate ${r(x, y) = r_{\phi}(x, y) -\beta (\log \pi_{\theta}(y\mid x) - \log \pi_\text{ref}(y\mid x))}$, which is then maximized using PPO~\cite{schulman2017proximal}.

\section{Direct Preference Optimization}\label{sec:DPO}

Recognizing the difficulties of scaling reinforcement learning to large language models, our aim is to derive a streamlined policy optimization method that works directly with preference data. In contrast to conventional RLHF approaches that learn a reward and then perform RL, our methodology uses a specific reward model parameterization allowing us to compute the optimal policy analytically, sidestepping any RL optimization loop. As detailed below, our principal observation is that with a known mapping from rewards to optimal policies, one can reformulate the loss originally defined for rewards as a loss on the policy space itself. This transformation obviates the need for a separate reward model, yet it still honors well-established preference models such as the Bradley-Terry formulation. Consequently, the policy neural network encapsulates both the generative model and its implicit reward signal.

\textbf{Derivation of the DPO objective.} Our starting point is the same RL objective under a general reward function $r$ as considered in previous studies, specifically Eq.~\ref{eq:RL}. Prior analyses~\citep{peters2007reinforcement, peng2019advantage, korbak2022reinforcement, go2023aligning} demonstrate that the solution maximizing the KL-constrained reward objective in Eq.~\ref{eq:RL} is given by:

\begin{equation}\label{eq:op_policy}
    \pi_r(y\mid x) = \frac{1}{Z(x)}\pi_\text{ref}(y\mid x)\exp\left(\frac{1}{\beta}r(x, y)\right),
\end{equation}

where $Z(x) =\sum_{y}\pi_\text{ref}(y\mid x)\exp\left(\frac{1}{\beta}r(x, y)\right)$ serves as the normalization constant (partition function). Detailed steps are included in Appendix \ref{app:derivation1}. Even with an MLE estimate $r_{\phi}$ for the ground-truth reward function $r^*$, evaluating $Z(x)$ remains computationally burdensome~\citep{korbak2022reinforcement, go2023aligning}, limiting the practicality of this policy representation. To address this, one can rearrange Eq.~\ref{eq:op_policy} and instead express the reward as a function of its optimal policy $\pi_r$, the baseline $\pi_\text{ref}$, and the unobserved partition function $Z(\cdot)$. Specifically, after applying the logarithm to both sides of Eq.~\ref{eq:op_policy} and performing algebraic manipulations, we arrive at:

\begin{equation}\label{eq:main_eq}
    r(x,y) =\beta \log \frac{\pi_r(y\mid x)}{\pi_\text{ref}(y\mid x)} + \beta \log Z(x).
\end{equation}

This reformulation applies to the true reward $r^*$ as well as its associated optimal policy $\pi^*$. Importantly, the Bradley-Terry preference model relies exclusively on reward differences between two choices, specifically, ${p^*(y_1 \succ y_2 \mid x) = \sigma(r^*(x, y_1) - r^*(x, y_2))}$. Substituting the above expression for $r^*(x,y)$ from Eq.~\ref{eq:main_eq} into Eq.~\ref{eq:bradley-terry}, the terms involving the partition function cancel. As a result, the human preference likelihood depends only on the optimal policy $\pi^*$ and the reference policy $\pi_\text{ref}$:

\begin{equation}\label{eq:objective}
    p^*(y_1\succ y_2 \mid x)=\frac{1}{1 + \exp\left(\beta \log \frac{\pi^*(y_2\mid x)}{\pi_\text{ref}(y_2\mid x)} - \beta \log \frac{\pi^*(y_1\mid x)}{\pi_\text{ref}(y_1\mid x)}\right)}
\end{equation}

This derivation can be found in Appendix~\ref{app:derivation2}. Although Eq.~\ref{eq:objective} utilizes the Bradley-Terry model, analogous derivations can be established for broader Plackett-Luce models~\citep{plackett1975analysis, luce2012individual}, as detailed in Appendix~\ref{app:plackett_luce_models}.

Given that the human preference distribution can be recast in terms of the optimal policy instead of the reward function, we can directly construct a maximum likelihood training objective for a policy parametrized by $\pi_\theta$. Mirroring the reward-based approach (see Eq.~\ref{eq:reward_model}), this leads to the policy learning objective:

\begin{equation}\label{eq:optimum_model}
    \mathcal{L}_\text{DPO}(\pi_{\theta}; \pi_\text{ref}) = -\mathbb{E}_{(x, y_w, y_l)\sim \mathcal{D}}\left[\log \sigma \left(\beta \log \frac{\pi_{\theta}(y_w\mid x)}{\pi_\text{ref}(y_w\mid x)} - \beta \log \frac{\pi_{\theta}(y_l\mid x)}{\pi_\text{ref}(y_l\mid x)}\right)\

In this framework, we approximate an implicit reward through an alternative form of parameterization, in which the optimal policy directly corresponds to $\pi_\theta$. Additionally, because our approach parallels the estimation of a reparametrized Bradley-Terry model, it inherits certain theoretical guarantees, including consistency, provided the distribution of preference data meets appropriate criteria \cite{bong2022generalized}. We explore additional theoretical implications of DPO and its connections to related literature in Section~\ref{sec:theory}.

\textbf{How does the DPO update operate?} To gain mechanistic insight into DPO, it is informative to consider the gradient of its objective function $\mathcal{L}_\text{DPO}$. The derivative with respect to $\theta$ is expressed as:

\begin{multline*}\label{eq:gradient}
    \nabla_\theta \mathcal{L}_\text{DPO}(\pi_\theta;\pi_\text{ref}) = \\ -\beta\mathbb{E}_{(x, y_w, y_l) \sim \mathcal{D}} \bigg[\underbrace{\sigma(\hat{r}_\theta(x, y_l) - \hat{r}_\theta (x, y_w))}_\text{assigns greater emphasis when the reward ordering is incorrect}\bigg[\underbrace{\nabla_\theta\log \pi(y_w \mid x)}_\text{promote $y_w$'s likelihood} - \underbrace{\nabla_\theta\log\pi(y_l \mid x)}_\text{reduce $y_l$'s likelihood}\bigg]\bigg],
\end{multline*}

Here, the implicit reward function is defined as $\hat{r}_\theta(x, y) = \beta \log \frac{\pi_\theta(y \mid x)}{\pi_\text{ref}(y \mid x)}$, reflecting the comparative preferences encoded by $\pi_\theta$ relative to $\pi_\text{ref}$ (see Section~\ref{sec:theory} for further discussion). Intuitively, this gradient step modifies $\mathcal{L}_\text{DPO}$ to favor higher probabilities for the chosen completions $y_w$, while reducing the probabilities of those less favored, $y_l$. The influence of each data point is governed by the extent to which the model's implicit reward misranks the completions, moderated by $\beta$, reflecting the strength of the imposed KL penalty. Our empirical results indicate this weighting mechanism is crucial—an unweighted or naive approach can lead the language model to produce degenerate outputs (see Appendix Table~\ref{tab:unlikelihood_generations} for details).

\textbf{DPO overview.} The standard DPO framework proceeds in two main steps: First, for each prompt $x$, two completions $y_1, y_2 \sim \pi_\text{ref}(\cdot \mid x)$ are sampled and annotated with human preference signals, yielding an offline preference dataset $\mathcal{D} = \{x^{(i)}, y_w^{(i)}, y_l^{(i)}\}_{i=1}^N$. Second, the language model $\pi_\theta$ is trained to minimize the loss $\mathcal{L}_\text{DPO}$, using the specified reference model $\pi_\text{ref}$, dataset $\mathcal{D}$, and a selected value of $\beta$. 
In real-world scenarios, it is preferable to leverage existing public preference datasets instead of producing new samples and collecting fresh human feedback. Since these datasets are commonly generated using a specific model, $\pi^\text{SFT}$, we choose $\pi_\text{ref} = \pi^\text{SFT}$ whenever this model is available. In situations where $\pi^\text{SFT}$ is not present, we set $\pi_\text{ref}$ by fitting it to the preferred completions $(x, y_w)$ through maximum likelihood estimation: ${\pi_\text{ref} = \argmax_{\pi}\mathbb{E}_{x, y_w \sim \mathcal{D}}\left[\log \pi(y_w \mid x)\right]}$. This approach reduces the mismatch between the unknown true reference distribution and the $\pi_\text{ref}$ utilized by DPO. Additional implementation details and hyperparameter choices can be found in Appendix~\ref{app:implementation}.

\section{Theoretical Analysis of DPO} This section provides a deeper theoretical understanding of the DPO method, elucidates its theoretical justifications, and draws connections between the strengths of DPO and the known shortcomings of actor-critic methods frequently applied in RLHF settings, such as PPO~\cite{schulman2017proximal}.

\label{sec:theory}

\subsection{Your Language Model Is Secretly a Reward Model} DPO eliminates the need for separately fitting a reward function and subsequently applying reinforcement learning by instead employing a single maximum likelihood objective. Notably, the loss function in Eq. \ref{eq:main_eq} aligns with the Bradley-Terry model when the reward is parameterized as $r^*(x, y) = \beta \log\frac{\pi^*_\theta(y \mid x)}{\pi_\text{ref}(y \mid x)}$. The model $\pi_{\theta}$ is thus optimized in a manner equivalent to the reward model objective in Eq. \ref{eq:reward_model}, given a suitable variable transformation. In this part, we formalize the theoretical foundation behind this parameterization, demonstrate that it does not limit the expressiveness of reward models that can be learned, and prove that it enables exact recovery of the optimal policy. We begin by formalizing an equivalence relation for reward functions.

\begin{definition}
Two reward functions $r(x, y)$ and $r'(x, y)$ are defined to be equivalent if and only if there exists a function $f$ such that ${r(x, y)-r'(x, y) = f(x)}$.
\end{definition}

This equivalence is straightforward to verify, and it naturally induces a partition over the space of reward functions into equivalence classes. The following lemmas clarify important implications of this equivalence:

\begin{lemma}\label{lemma:same_prefrence} 
Within the Plackett-Luce framework, which subsumes the Bradley-Terry model, any pair of reward functions that belong to the same equivalence class generate identical preference distributions.
\end{lemma}

\begin{lemma}\label{lemma:same_policy}
Reward functions that are equivalent in this sense lead to the same optimal solution for the constrained reinforcement learning problem.
\end{lemma}

The arguments supporting these lemmas are elementary; detailed proofs can be found in Appendix~\ref{app:lemma1}. The first lemma encapsulates a well-documented indeterminacy in the Plackett-Luce models \cite{plackett1975analysis}, highlighting that preference data alone cannot distinguish among certain reward transformations. To circumvent this, extra identifiability criteria are generally introduced when estimating rewards via maximum likelihood, as noted in Eq.~\ref{eq:reward_model} and discussed in \cite{bong2022generalized}. The second lemma establishes that all members of an equivalence class of reward functions result in the same optimal policy; consequently, our goal may be to recover any single representative of the optimal class. The subsequent theorem, with proof provided in Appendix~\ref{app:thm1}, formalizes this insight:

\begin{theorem}\label{thm:main}
    Under mild conditions, every reward class that is consistent with the Plackett-Luce model (including Bradley-Terry as a special case) admits a representation in the form ${r(x, y) = \beta \log \frac{\pi(y\mid x)}{\pi_\text{ref}(y\mid x)}}$ for an appropriately chosen model $\pi(y\mid x)$ and fixed reference model $\pi_\text{ref}(y\mid x)$.
\end{theorem}

\begin{sproof}
    Take any reward function $r(x, y)$, which defines its corresponding optimal policy $\pi_r(y \mid x)$ as specified by Eq.~\ref{eq:op_policy}. We demonstrate that within the equivalence class of $r$, there is always a representative of the form indicated above. Define the mapping $f$ as
\begin{equation}
    f(r; \pi_\text{ref}, \beta)(x, y) = r(x, y) - \beta\log\sum_{y}\pi_\text{ref}(y\mid x)\exp\left(\frac{1}{\beta}r(x, y)\right)
\end{equation}
This operation $f$ acts by subtracting the log partition function—dependent solely on the context $x$—from $r(x, y)$. As the normalization term varies only with $x$, $f(r; \pi_\text{ref}, \beta)(x, y)$ stays within the equivalence class defined by $r(x, y)$. By inserting for $r$ the right side of Eq.~\ref{eq:main_eq} (which is valid for any reward function), we obtain $f(r; \pi_\text{ref}, \beta)(x, y) = \beta \log \frac{\pi_r(y\mid x)}{\pi_\text{ref}(y\mid x)}$. Therefore, $f$ produces an element of the equivalence class of $r$ that has the desired structure, showing that the reparameterization retains full generality for modeling rewards within the Plackett-Luce family.
\end{sproof}

Theorem~\ref{thm:main} can also be interpreted as identifying the specific reward function selected by the DPO reparameterization within each equivalence class; namely, it selects the reward function that fulfills the following condition:

\begin{equation}\label{eq:lag_p}
     \sum_{y}\underbrace{\pi_\text{ref}(y\mid x)\exp\left(\frac{1}{\beta}r(x, y)\right)}_{=\pi(y\mid x)\text{, using Thm.~\ref{thm:main} reparam.}} = 1,
\end{equation}

which means that $\pi(y\mid x)$ forms a legitimate probability distribution (all entries are non-negative and total to one). When considering Eq.~\ref{eq:op_policy}, it becomes clear that Eq.~\ref{eq:lag_p} represents the partition function associated with the optimal policy produced by the reward $r(x, y)$. The essential idea behind the DPO algorithm is that one can place additional constraints on the otherwise underdetermined Plackett-Luce (and, specifically, Bradley-Terry) family of preference models so that, while the expressiveness of the class of reward models remains unaffected, the resulting optimal policy from Eq.~\ref{eq:op_policy} becomes computationally manageable for every prompt $x$.

\subsection{Instability of Actor-Critic Algorithms} Our framework further serves to explain sources of instability present in commonly used actor-critic approaches for RLHF, such as PPO. We adhere to the RLHF process and, in particular, the RL fine-tuning phase as detailed in Section \ref{section:prelims}. There exists a relationship to the control as inference paradigm \cite{levine2018reinforcement}, applicable to the constrained RL scenario described in \ref{eq:RL}. By assuming a parameterized policy $\pi_{\theta}(y\mid x)$, the objective becomes minimizing $\mathbb{D}_{\text{KL}}[\pi_{\theta}(y|x) \mid \mid \pi^*(y\mid x)]$, with $\pi^*$ representing the optimal policy per Eq.~\ref{eq:optimum_model}, determined by the reward $r_{\phi}(y, x)$. Through manipulation, this yields the following objective function:

\begin{equation}\label{eq:AC}
    \max_{\pi_{\theta}}\mathbb{E}_{\pi_{\theta}(y\mid x)}\bigg[\underbrace{r_{\phi}(x, y) -\beta\log\sum_{y}\pi_\text{ref}(y\mid x)\exp\left(\frac{1}{\beta}r_{\phi}(x, y)\right)}_{f(r_{\phi}, \pi_\text{ref}, \beta)} - \underbrace{\beta\log\frac{\pi_{\theta}(y\mid x)}{\pi_\text{ref}(y\mid x)}}_{\text{KL}}\bigg]
\end{equation}

The objective here aligns with that optimized in earlier studies \citep{ziegler2020finetuning, stiennon2022learning, bai2022training, ouyang2022training}, which also utilize a DPO-equivalent reward for the reward family $r_{\phi}$. In this context, the normalization component within $f(r_{\phi}, \pi_\text{ref}, \beta)$ can be viewed as the reference policy $\pi_\text{ref}$'s soft value function. Although this term is not required for obtaining the optimal policy, omitting it can cause the policy gradient estimator to exhibit large variance, thereby potentially destabilizing the learning process. One option is to use a learned value function to approximate the normalization term; however, this itself can pose significant optimization challenges. Prior studies have addressed normalization by referencing a human completion baseline—effectively employing a single-sample Monte Carlo estimate for the normalization. By contrast, the DPO reparameterization produces a reward formulation that eliminates the necessity for any baseline.

\section{Experiments}
We proceed to assess the effectiveness of DPO in training policies based solely on preference data. Our first investigation uses a controlled text generation scenario to examine DPO's efficiency in balancing reward maximization against KL-divergence minimization relative to the reference policy, benchmarking it against established preference learning methods such as PPO. Subsequently, we expand our analysis to larger-scale models and more complex RLHF tasks, including applications in summarization and dialogue. Our findings indicate that DPO achieves competitive, and often superior, performance compared to strong baselines like PPO-based RLHF, as well as against strategies that select the top result from $N$ trajectories using a learned reward model—remarkably, this holds with minimal hyperparameter adjustment. Prior to discussing empirical outcomes, we detail our experimental protocols; further specifics can be found in Appendix~\ref{app:exp_details}.

\textbf{Tasks.} We evaluate our methods across three open-ended text generation settings. In all cases, algorithms learn a policy using a preference dataset $\mathcal{D}=\bigl\{x^{(i)}, y_w^{(i)}, y_l^{(i)}\bigr\}_{i=1}^N$. The first task, \textbf{controlled sentiment generation}, uses $x$ as the beginning segment of a movie review sampled from the IMDb dataset \cite{maas-EtAl:2011:ACL-HLT2011}, with the objective that the policy generates a continuation $y$ expressing positive sentiment. For systematic assessment, we synthetically produce preference pairs based on generations scored by a pre-trained sentiment classifier, ensuring $p(\text{positive}\mid x,y_w)>p(\text{positive}\mid x,y_l)$. For SFT, we train GPT-2-large via fine-tuning on the IMDB train split, continuing until convergence (see App~\ref{app:sentiment_details} for additional information). The second task, \textbf{summarization}, sets $x$ as a Reddit forum post, with the goal for the policy to output a summary $y$ capturing the main ideas. As in previous research, we utilize the Reddit TL;DR dataset \citep{volske-etal-2017-tl}, combined with human-labeled preferences provided by \citeauthor{stiennon2022learning}. We employ a supervised fine-tuned (SFT) model trained on summaries written by humans\footnote{\url{https://huggingface.co/CarperAI/openai_summarize_tldr_sft}} and perform RLHF using the TRLX \citep{leandro_von_werra_2023_7790115} framework. The preference data were originally gathered by \citeauthor{stiennon2022learning} on outputs from a similar but distinct SFT model. The third task, \textbf{single-turn dialogue}, involves $x$ as a user prompt, which may vary from scientific inquiries to personal advice requests. Here, the policy's responsibility is to deliver an informative and engaging answer $y$. For this, we rely on the Anthropic Helpful and Harmless dialogue dataset \citep{bai2022training}, containing 170k interactions between users and a language model assistant. Each dialogue concludes with two model-generated responses, annotated with human preferences indicating which was favored. In this scenario, lacking a ready SFT model, we instead fine-tune a publicly available language model exclusively on responses marked as preferred to establish our SFT baseline.

\textbf{Evaluation.} Our experimental assessment employs two complementary evaluation strategies. To rigorously measure how well each algorithm achieves the goal of constrained reward maximization, we first focus on controlled sentiment generation. Here, for every method, we characterize performance via the trade-off frontier between the attained reward and the KL-divergence relative to a reference policy. Because access to an exact reward function (specifically, a sentiment classifier) is available, this frontier is directly computable. By contrast, in real-world scenarios where the reward function remains unobservable, we assess algorithms in terms of their \textit{win rate} against a baseline policy. For the tasks of summarization and single-turn dialogue, we rely on GPT-4 as a surrogate for human judges, evaluating output quality and response helpfulness. For summarization tasks, the provided test set reference summaries serve as the baseline; in the dialogue domain, the test set’s preferred responses are the benchmark. Although prior work indicates that language models can outperform conventional automated evaluation metrics \citep{Chen2023ExploringTU}, we reinforce our reliance on GPT-4 as an evaluator by conducting a human-subject study, the findings of which are presented in Sec.~\ref{sec:human-judgments}. The results reveal strong concordance between GPT-4 assessments and human judgments, with the alignment between GPT-4 and individual humans typically matching or even surpassing agreement among different human annotators.

\textbf{Methods.} Alongside DPO, we benchmark a selection of established techniques for adapting language models to align with human preferences. For baseline prompting, we implement zero-shot inference with \textbf{GPT-J} \citep{gpt-j} in summarization, and two-shot prompting with \textbf{Pythia-2.8B} \citep{biderman2023pythia} in dialogue tasks. Additional comparisons include the \textbf{SFT} model, as well as \textbf{Preferred-FT}, where supervised fine-tuning targets the selected completion $y_w$—derived either from the SFT model (in controlled sentiment generation and summarization) or a generic language model (for single-turn dialogue). We further evaluate the \textbf{Unlikelihood}~\citep{welleck2019neural} approach, which updates model parameters to both maximize likelihood for $y_w$ and actively decrease the likelihood of $y_l$; here, the unlikelihood penalty term is weighted by an adjustable coefficient $\alpha\in[0,1]$. We also include \textbf{PPO} \citep{schulman2017proximal}, in which reward models are trained from preference annotations, as well as \textbf{PPO-GT}, an oracle variant that utilizes ground truth rewards—possible in controlled sentiment scenarios. For the sentiment tasks, we employ two PPO-GT implementations: a standard out-of-the-box solution \cite{leandro_von_werra_2023_7790115}, and a customized version which includes reward normalization and hyperparameter adjustments for enhanced results (these adjustments also apply when running PPO with learned rewards). Lastly, we include a \textbf{Best of $N$} baseline that, for each query, generates $N$ candidate responses from the SFT model (or Preferred-FT in dialogue tasks) and selects the candidate receiving the top score according to the learned reward model. Although this technique can yield strong performance by separating reward model accuracy from PPO dynamics, it remains computationally expensive, since evaluating $N$ completions per query is required during testing even for moderate $N$.

\subsection{How well can DPO optimize the RLHF objective?}

Standard RLHF methods utilize a KL-constrained reward maximization framework, which seeks to maximize reward while simultaneously limiting the policy's divergence from a given reference policy. As such, a proper assessment of algorithms necessitates consideration of both achieved reward and KL divergence; maximizing reward at the cost of a disproportionately large KL increase may not be optimal. In Figure~\ref{fig:frontier-tldr-main}, the tradeoff curve between reward and KL for several algorithms is depicted in the sentiment task context. For each method, we conduct several independent training sessions, varying the policy conservativeness hyperparameter in each trial (using target KL values $\in\{3,6,9,12\}$ for PPO, $\beta \in \{0.05,0.1,1,5\}$, and $\alpha\in\{0.05,0.1,0.5,1\}$ for the unlikelihood approach, while randomizing seeds for preferred-FT), totaling 22 experiments. At intervals of 100 training steps, up to convergence, each trained policy is evaluated using a test set of prompts, with both mean reward (under the true reward metric) and average sequence-level KL\footnote{Defined as the cumulative sum of KL-divergences at each timestep.} relative to the reference policy $\text{KL}\left(\pi\mid \mid \pi_\text{ref}\right)$ computed. The results indicate that DPO establishes a markedly more efficient reward-KL frontier compared to the alternatives, achieving high rewards at comparatively low KL values. This outcome is noteworthy for several reasons. Most notably, while both DPO and PPO aim to optimize an identical objective, DPO does so with superior efficiency—the reward-KL tradeoff realized by DPO outperforms PPO consistently. Additionally, DPO’s performance remains better than that of PPO even when PPO operates with access to ground-truth reward signals (PPO-GT).

\subsection{Can DPO scale to real preference datasets?}
\label{sec:dpo-real-datasets}
We proceed to assess the ability of DPO to scale by examining its fine-tuning performance on real-world tasks, namely summarization and single-turn dialogue. In the context of summarization, it is known that widely used automated metrics like ROUGE often do not align well with human judgments~\citep{stiennon2022learning}. Previous research indicates that leveraging human feedback to fine-tune language models with PPO leads to improvements in summary quality. To compare various approaches, we generate completions on the test partition of the TL;DR summarization dataset, measuring the average win rate of these outputs against the reference completions provided in the test set. For each method, completions are generated across a range of temperatures from 0.0 to 1.0, and the associated win rates are depicted in Figure~\ref{fig:frontier-tldr-main} (right). DPO, PPO, and Preferred-FT all share the same GPT-J SFT model for fine-tuning\footnote{\url{https://huggingface.co/CarperAI/openai_summarize_tldr_sft}}. Our findings indicate that at a sampling temperature of 0.0, DPO attains a win rate of about 61\%, surpassing PPO, which achieves roughly 57\% at its optimal temperature. Furthermore, DPO attains a superior peak win rate relative to the $N$ baseline. It is important to mention that we did not extensively search over the DPO $\beta$ hyperparameter, and thus the reported results might not reflect the upper limit of DPO's capability. Additionally, DPO demonstrates greater resilience to changes in sampling temperature compared to PPO, whose performance can deteriorate to the level of the unadapted GPT-J at higher temperatures. The Preferred-FT method, in contrast, yields negligible gains over the SFT baseline. Direct human evaluation of DPO versus PPO, discussed in Section~\ref{sec:human-judgments}, further shows that DPO samples at a temperature of 0.25 were selected in 58\% of cases when compared with PPO samples generated at 0.

For the single-turn dialogue setting, we conduct experiments on a filtered portion of the Anthropic HH dataset’s test split \citep{bai2022training}, specifically focusing on samples containing a single exchange between a human and assistant. For the GPT-4-based assessments, we measure the win rate of each method by referencing the preferred completions present in the test set. In the absence of a standard SFT model for this benchmark, our procedure begins with the Pythia-2.8B model pre-trained weights; we then apply Preferred-FT to fine-tune a reference model on the preferred completions, aligning output distributions accordingly, before further optimization with DPO. Additionally, we benchmark performance against two alternatives: the optimal output among 128 Preferred-FT completions (as empirical results indicate the Best of $N$ approach plateaus at 128, see Appendix Figure~\ref{fig:best-of-n}), and a two-shot prompting strategy applied to the original Pythia-2.8B model. Across these methods, DPO matches or exceeds the best-performing results when appropriate temperatures are selected. We further include an RLHF model trained with PPO on the Anthropic HH dataset \footnote{\url{https://huggingface.co/reciprocate/ppo_hh_pythia-6B}}, developed by a reputable third party \footnote{\url{https://github.com/CarperAI/trlx/tree/main/examples/hh}}; however, we are unable to identify any prompting or temperature strategy that achieves results surpassing those of the original Pythia-2.8B base model. Drawing on insights from the TL;DR task and the observation that both approaches optimize the same reward objective, we treat Best of 128 as an approximate indicator of PPO model performance. Summing up, DPO emerges as the sole approach that not only delivers significant efficiency gains but also consistently outperforms the preferred completions within the Anthropic HH dataset, matching or exceeding the effectiveness of the computationally costly Best of 128 reference. Figure~\ref{fig:dialogue-main} further illustrates that DPO attains its peak performance after relatively few training steps.

\subsection{Generalization to a new input distribution}

To further investigate how PPO and DPO respond to changes in data distribution, we assess the policies derived from our Reddit TL;DR summarization study on an out-of-distribution dataset: news articles from the CNN/DailyMail test split \citep{nallapati-etal-2016-abstractive}. We employ the optimal sampling temperatures from the TL;DR experiments (0 and 0.25). The findings are summarized in Table~\ref{tab:ood}. Evaluation involves calculating the win rate for GPT-4 versus the human-authored reference summaries, utilizing the same GPT-4 (C) prompt from the TL;DR setting but substituting ``forum post'' with ``news article.'' On this alternative dataset, DPO substantially outperforms PPO. These outcomes suggest that DPO demonstrates comparable, if not superior, generalization abilities relative to PPO, despite DPO not utilizing the extra unlabeled Reddit TL;DR prompts available to PPO.

\subsection{Validating GPT-4 judgments with human judgments}
\label{sec:human-judgments}
To assess the dependability of GPT-4's scoring, we conduct a human evaluation based on results from the TL;DR summarization task and two distinct GPT-4 prompts. The \textbf{GPT-4 (S)} (simple) prompt instructs GPT-4 to choose the summary that better covers the key information in a post. Conversely, the \textbf{GPT-4 (C)} (concise) prompt includes an explicit instruction to consider conciseness; we consider this prompt as GPT-4 using the \textbf{GPT-4 (S)} prompt often favors summaries that are lengthier and more redundant compared to human preferences. Full prompt text is provided in Appendix~\ref{app:prompts}. We perform three evaluation pairings—DPO at temperature 0.25 (highest-performing), PPO at temperature 1.0 (lowest-performing), and SFT at temperature 0.25 (intermediate)—selected to capture a spectrum of generation quality. Each of these models is compared against PPO with greedy decoding (its best-performing temperature). Our findings indicate that GPT-4, with both prompt variants, aligns with human raters at a rate similar to human-human agreement, suggesting that GPT-4 serves as a credible substitute for human judgment (for comparisons involving DPO and PPO-1, we obtain multiple human ratings, constrained by the number of available raters). Overall, results from the \textbf{GPT-4 (C)} prompt most closely mirror those from human raters, which motivates our use of this prompt for key analyses in Section~\ref{sec:dpo-real-datasets}. More details about the human evaluation, including the interface and rater information, are provided in Appendix~\ref{app:human-study}.

\section{Discussion}
Preference-based learning offers an effective and scalable approach for developing capable and aligned language models. In this work, we presented DPO, a straightforward method for training language models directly from preference data, eliminating the need for reinforcement learning. Instead of reframing the preference learning problem into a conventional reinforcement learning (RL) framework to take advantage of existing RL algorithms, DPO establishes a correspondence between policies of language models and reward functions. This enables the direct optimization of language models to align with human preferences using a basic cross-entropy loss, foregoing both reinforcement learning procedures and any compromise on generality. With minimal hyperparameter tuning, DPO achieves performance on par with or superior to contemporary RLHF methods, such as those leveraging PPO. As a result, DPO substantially lowers the barriers for using preference data to train language models.

\textbf{Limitations \& Future Work.} This study highlights several open directions for subsequent investigation. A central question concerns the out-of-distribution generalization capabilities of DPO policies as compared to models trained on explicit reward signals. Preliminary evidence points to DPO exhibiting generalization behavior akin to that of PPO-trained models, but a more in-depth analysis is required. For instance, the potential for self-labeled DPO training to utilize unlabeled prompts effectively remains to be thoroughly explored. Another topic of interest is understanding how phenomena like reward over-optimization arise in direct preference optimization, and whether the slight reduction in performance observed in Figure~\ref{fig:dialogue-main}-right provides an example of this issue. Furthermore, while our experiments include models up to 6B parameters, scaling DPO to much larger, cutting-edge models represents a promising avenue. In the realm of evaluation, we note that the win rates assigned by GPT-4 are sensitive to the prompting technique; future studies could investigate more robust strategies for obtaining high-fidelity assessments from automated evaluators. Beyond language model training, there is also significant potential for extending DPO to train generative models across other modalities.